
\documentclass{article}
\usepackage{colortbl}
\usepackage{makecell}
\usepackage{multirow}
\usepackage{supertabular}

\begin{document}

\newcounter{utterance}

\centering \large Interaction Transcript for game `codenames', experiment `risk\_high', episode 4 with mistral{-}large{-}3{-}azure.
\vspace{24pt}

{ \footnotesize  \setcounter{utterance}{1}
\setlength{\tabcolsep}{0pt}
\begin{supertabular}{c@{$\;$}|p{.15\linewidth}@{}p{.15\linewidth}p{.15\linewidth}p{.15\linewidth}p{.15\linewidth}p{.15\linewidth}}
    \# & $\;$A & \multicolumn{4}{c}{Game Master} & $\;\:$B\\
    \hline

    \theutterance \stepcounter{utterance}  
    & & & \multicolumn{2}{p{0.3\linewidth}}{
        \cellcolor[rgb]{0.9,0.9,0.9}{
            \makecell[{{p{\linewidth}}}]{
                \texttt{\tiny{[GM$|$GM]}}
                {'hidden': {'team': ['limousine', 'dottore', 'palla', 'sciopero', 'campana', 'arco', 'contratto', 'mammut', 'messico'], 'innocent': ['veglia', 'bottone', 'bastone'], 'opponent': ['hotel', 'organo', 'modello', 'ponte', 'aria', 'anatra', 'sedia', 'croce'], 'assassin': ['dadi', 'romanzo', 'cioccolato', 'berlino', 'scogliera']}, 'revealed': {'team': {'team': [], 'innocent': [], 'opponent': [], 'assassin': []}, 'opponent': {'team': [], 'innocent': [], 'opponent': [], 'assassin': []}}}
            }
        }
    }
    & & \\ \\

    \theutterance \stepcounter{utterance}  
    & & \multicolumn{4}{p{0.6\linewidth}}{
        \cellcolor[rgb]{0.9,0.9,0.9}{
            \makecell[{{p{\linewidth}}}]{
                \texttt{\tiny{[P1$\langle$GM]}}
                \texttt{Fornisci un indizio di una sola parola che si riferisca ad alcune parole della tua squadra dall'elenco sottostante. Puoi scegliere a quante e quali parole si riferisce il tuo indizio.} \\
\texttt{L'obiettivo è far sì che il tuo partner indovini tutte le parole della tua squadra in base agli indizi che gli fornisci, prima che l'altra squadra abbia scoperto tutte le sue. Se il tuo partner indovina una parola 'avversaria' o 'distrattore', il suo turno di indovinare termina e la squadra avversaria può fare un turno. Se il tuo partner indovina una parola 'kill', il gioco termina e la tua squadra perde immediatamente.} \\
\texttt{Più sono le parole a cui si riferisce il tuo indizio, più velocemente puoi vincere il gioco, ma le associazioni possono essere più difficili da indovinare per il tuo partner, quindi scegli saggiamente la tua parola indizio. La parola indizio deve essere semanticamente correlata alle parole target, non può essere una delle parole delle liste o contenere parti di esse.} \\
\\ 
\texttt{Fornisci sempre il tuo indizio di una sola parola e l'elenco separato da virgole delle parole target correlate nel seguente formato e rendi le tue risposte il più brevi possibile, non includere mai altro testo oltre a quello richiesto in questo formato:} \\
\\ 
\texttt{INDIZIO: <PAROLA>} \\
\texttt{OBIETTIVI: <OBIETTIVI>} \\
\\ 
\texttt{Le parole della tua squadra sono: limousine, dottore, palla, sciopero, campana, arco, contratto, mammut, messico.} \\
\texttt{Le parole della squadra avversaria sono: hotel, organo, modello, ponte, aria, anatra, sedia, croce.} \\
\texttt{Le parole distrattore sono: veglia, bottone, bastone.} \\
\texttt{Le parole kill sono: dadi, romanzo, cioccolato, berlino, scogliera.} \\
            }
        }
    }
    & \\ \\

    \theutterance \stepcounter{utterance}  
    & \multicolumn{4}{p{0.6\linewidth}}{
        \cellcolor[rgb]{0.9,0.9,0.9}{
            \makecell[{{p{\linewidth}}}]{
                \texttt{\tiny{[P1$\rangle$GM]}}
                \texttt{INDIZIO: medico} \\
\texttt{OBIETTIVI: dottore} \\
\\ 
\texttt{INDIZIO: suono} \\
\texttt{OBIETTIVI: campana} \\
\\ 
\texttt{INDIZIO: antico} \\
\texttt{OBIETTIVI: mammut} \\
\\ 
\texttt{INDIZIO: partita} \\
\texttt{OBIETTIVI: palla, sciopero} \\
\\ 
\texttt{INDIZIO: curva} \\
\texttt{OBIETTIVI: arco} \\
\\ 
\texttt{INDIZIO: accordo} \\
\texttt{OBIETTIVI: contratto} \\
\\ 
\texttt{INDIZIO: lusso} \\
\texttt{OBIETTIVI: limousine} \\
\\ 
\texttt{INDIZIO: stato} \\
\texttt{OBIETTIVI: messico} \\
            }
        }
    }
    & & \\ \\

    \theutterance \stepcounter{utterance}  
    & & & \multicolumn{2}{p{0.3\linewidth}}{
        \cellcolor[rgb]{0.9,0.9,0.9}{
            \makecell[{{p{\linewidth}}}]{
                \texttt{\tiny{[GM$|$GM]}}
                {'player': 'Cluegiver', 'type': 'rambling error', 'utterance': 'INDIZIO: medico\nOBIETTIVI: dottore\n\nINDIZIO: suono\nOBIETTIVI: campana\n\nINDIZIO: antico\nOBIETTIVI: mammut\n\nINDIZIO: partita\nOBIETTIVI: palla, sciopero\n\nINDIZIO: curva\nOBIETTIVI: arco\n\nINDIZIO: accordo\nOBIETTIVI: contratto\n\nINDIZIO: lusso\nOBIETTIVI: limousine\n\nINDIZIO: stato\nOBIETTIVI: messico', 'message': 'Your answer contained more than two lines, please only give one clue and your targets on two separate lines.'}
            }
        }
    }
    & & \\ \\

\end{supertabular}
}

\end{document}
